\documentclass[a4paper,oneside,11pt]{book}

% Package dasar
\usepackage[utf8]{inputenc}     % untuk karakter UTF-8
\usepackage[T1]{fontenc}        % encoding font
\usepackage{graphicx}
% Manual redefinition untuk bahasa Indonesia (pengganti babel)
\renewcommand{\chaptername}{Bab}
\renewcommand{\contentsname}{Daftar Isi}
\renewcommand{\listfigurename}{Daftar Gambar}
\renewcommand{\figurename}{Gambar}
\renewcommand{\tablename}{Tabel}
\renewcommand{\appendixname}{Lampiran}
\renewcommand{\bibname}{Daftar Pustaka}
\renewcommand{\indexname}{Indeks}
\providecommand{\abstractname}{Abstrak}
\providecommand{\chaptertitlename}{\chaptername}

% Package tambahan untuk layout dan formatting
\usepackage{geometry}           % untuk margin
\usepackage{setspace}           % untuk spasi baris
\usepackage{fancyhdr}           % untuk header/footer
\usepackage{hyperref}           % untuk link dan bookmark
\usepackage{listings}           % untuk code
\usepackage{xcolor}             % untuk warna
\usepackage{float}              % untuk posisi gambar
\usepackage{caption}            % untuk caption gambar
\captionsetup{hypcap=false}

% Pengaturan layout
\geometry{left=2cm, right=2cm, top=2cm, bottom=2cm}
\onehalfspacing  % spasi 1.5

% Menghindari warning fancyhdr headheight
\setlength{\headheight}{14pt}

\lstset{
    basicstyle=\ttfamily\small,
    keywordstyle=\color{blue}\bfseries,
    commentstyle=\color{gray},
    stringstyle=\color{red},
    identifierstyle=\color{black},
    numbers=left,
    numberstyle=\tiny\color{gray},
    numbersep=10pt,
    frame=single,
    backgroundcolor=\color{gray!10},
    breaklines=true,
    showstringspaces=false,
    tabsize=4,
    xleftmargin=2cm,
    xrightmargin=1cm,
    captionpos=b
}

% Pengaturan hyperref
\hypersetup{
    colorlinks=true,
    linkcolor=black,
    filecolor=magenta,
    urlcolor=blue,
    citecolor=red
}

% Info dokumen
\title{LAPORAN PRAKTIKUM \\
PENGUJIAN PERANGKAT LUNAK \\
LAPORAN 12 \\
BEHAVIOR DRIVEN DEVELOPMENT (BDD) DENGAN GHERKIN DAN CUCUMBER}
\author{Disusun oleh: Matthew Hayunaji Priantara\\24/536179/SV/24400}

\begin{document}

% Title Page
\begin{titlepage}
\begin{center}

\vspace{1cm}
\begin{Large}
\textbf{LAPORAN PRAKTIKUM} \\
\textbf{PENGUJIAN PERANGKAT LUNAK} \\
\textbf{PERTEMUAN 12} \\
\textbf{BEHAVIOR DRIVEN DEVELOPMENT (BDD) DENGAN GHERKIN DAN CUCUMBER} \\
\end{Large}
\vspace{2cm}

% HINT GAMBAR: Ganti logo UGM atau sesuaikan ukuran jika perlu
\includegraphics[height=8cm]{../lambang ugm.png}\\
\vspace{3cm}

\begin{Large}
\textbf{Disusun oleh:}\\
\textbf{Matthew Hayunaji Priantara}\\
\textbf{24/536179/SV/24400}\\
\end{Large}

\vfill
\begin{Large}
\textbf{PROGRAM STUDI D-IV TEKNOLOGI REKAYASA PERANGKAT LUNAK}\\
\textbf{DEPARTEMEN TEKNIK ELEKTRO DAN INFORMATIKA}\\
\textbf{SEKOLAH VOKASI}\\
\textbf{UNIVERSITAS GADJAH MADA}\\
\textbf{YOGYAKARTA}\\
\textbf{2026}\\
\end{Large}

\end{center}
\end{titlepage}

% Header/Footer
\pagestyle{fancy}
\fancyhf{}
\fancyhead[L]{Laporan Praktikum}
\fancyhead[R]{Pengujian Perangkat Lunak}
\fancyfoot[C]{\thepage}

% Daftar Isi
\tableofcontents
\newpage

% DAFTAR GAMBAR
\cleardoublepage
\addcontentsline{toc}{chapter}{Daftar Gambar}
\listoffigures
\newpage

% BAB 1: PENDAHULUAN
\chapter{Pendahuluan}
\section{Tujuan Praktikum}
\begin{enumerate}
  \item Memahami metodologi \textit{Behavior Driven Development} (BDD) dan penggunaan sintaks Gherkin dalam menyusun skenario pengujian berbasis perilaku pengguna.
  \item Mengimplementasikan pengujian otomatis menggunakan framework Cucumber yang terintegrasi dengan Selenium WebDriver dan JUnit 5 di lingkungan pemrograman Java.
  \item Memahami penerapan Cucumber Hooks (\texttt{@Before} dan \texttt{@After}) untuk mengelola siklus hidup pengujian secara terstruktur dan efisien.
  \item Menguasai teknik penyaringan skenario uji (\textit{test filtering}) menggunakan Tag Gherkin pada Test Runner JUnit 5 Suite untuk kebutuhan eksekusi yang spesifik.
\end{enumerate}

% BAB 2: HASIL DAN PEMBAHASAN
\chapter{Hasil dan Pembahasan}

\section{Deskripsi Pengujian}
Praktikum ini berfokus pada penerapan pengujian otomatis dengan pendekatan \textit{Behavior Driven Development} (BDD) menggunakan Cucumber. Kasus uji yang digunakan adalah proses autentikasi (login) pada aplikasi e-commerce SauceDemo. Pengujian dilakukan melalui skenario-skenario perilaku yang ditulis menggunakan bahasa natural terstruktur (Gherkin Syntax) dan diterjemahkan menjadi kode eksekusi Selenium WebDriver.

Secara garis besar, tahapan pengerjaan dalam praktikum ini meliputi:
\begin{enumerate}
    \item \textbf{Konfigurasi Dependensi BDD}: Menyiapkan lingkungan kerja Maven melalui berkas \texttt{pom.xml} dengan menambahkan pustaka JUnit 5, Cucumber, dan Selenium.
    \item \textbf{Tugas 1 (Implementasi BDD dan Cucumber)}: Membuat berkas fitur Gherkin (\texttt{login.feature}), menyusun kelas-kelas \textit{Page Object Model} (\texttt{basePage}, \texttt{loginPage}, dan \texttt{inventoryPage}), menghubungkannya melalui \texttt{loginStepDef}, serta mengonfigurasi JUnit 5 Suite Runner (\texttt{LoginSuiteRunner}).
    \item \textbf{Tugas 2 (Eksplorasi Tag Gherkin dan Filter Tag)}: Melakukan penyaringan pengujian menggunakan parameter konfigurasi sehingga JUnit Suite hanya mengeksekusi skenario negatif (\texttt{@Negative}).
    \item \textbf{Tugas 3 (Penerapan Saringan Test Spesifik)}: Melakukan penyaringan pengujian spesifik untuk pengguna standar (\texttt{@StandardUser}).
    \item \textbf{Poin Tambahan (Bonus)}: Menerapkan Cucumber Hooks untuk mengatur otomatisasi persiapan peramban (\textit{setup}) dan pembersihan sumber daya (\textit{cleanup}).
\end{enumerate}

\section{Konfigurasi Dependensi BDD (pom.xml)}
Sebelum menulis kode pengujian, konfigurasi dependensi didefinisikan pada berkas \texttt{pom.xml} proyek. Hal ini bertujuan untuk mengunduh pustaka-pustaka yang diperlukan dalam menyusun pengujian berbasis Cucumber, JUnit 5 Suite, dan Selenium.

\begin{lstlisting}[language=XML, caption=Konfigurasi Dependensi pada pom.xml]
<dependencies>
    <!-- JUnit Jupiter untuk Assertion -->
    <dependency>
        <groupId>org.junit.jupiter</groupId>
        <artifactId>junit-jupiter</artifactId>
        <version>5.13.4</version>
        <scope>test</scope>
    </dependency>
    
    <!-- Selenium Java untuk Otomasi Browser -->
    <dependency>
        <groupId>org.seleniumhq.selenium</groupId>
        <artifactId>selenium-java</artifactId>
        <version>4.43.0</version>
        <scope>compile</scope>
    </dependency>
    
    <!-- Cucumber Java untuk Sintaks Gherkin -->
    <dependency>
        <groupId>io.cucumber</groupId>
        <artifactId>cucumber-java</artifactId>
        <version>7.33.0</version>
        <scope>compile</scope>
    </dependency>
    
    <!-- Cucumber JUnit Platform Engine untuk Integrasi JUnit 5 -->
    <dependency>
        <groupId>io.cucumber</groupId>
        <artifactId>cucumber-junit-platform-engine</artifactId>
        <version>7.33.0</version>
        <scope>test</scope>
    </dependency>
    
    <!-- JUnit Platform Suite untuk Runner -->
    <dependency>
        <groupId>org.junit.platform</groupId>
        <artifactId>junit-platform-suite</artifactId>
        <version>1.13.4</version>
        <scope>test</scope>
    </dependency>
</dependencies>
\end{lstlisting}

\subsection{Analisis Dependensi}
Berdasarkan deklarasi dependensi di atas, fungsionalitas masing-masing pustaka dijabarkan sebagai berikut:
\begin{enumerate}
    \item \textbf{junit-jupiter}: Berperan sebagai fondasi pengujian unit dan integrasi. Pustaka ini menyediakan fungsi-fungsi penegasan (\textit{assertions}) seperti \texttt{assertEquals} dan \texttt{assertTrue} untuk memvalidasi apakah hasil pengujian aktual sesuai dengan skenario yang diharapkan.
    \item \textbf{selenium-java}: Pustaka utama WebDriver yang memungkinkan program untuk mengendalikan peramban web (seperti Google Chrome) secara otomatis. Pustaka ini digunakan untuk mencari elemen HTML, mensimulasikan input teks, klik tombol, serta bernavigasi antar halaman web.
    \item \textbf{cucumber-java}: Framework utama untuk menulis tes dengan gaya BDD. Pustaka ini berfungsi memetakan langkah-langkah skenario teks terstruktur (Gherkin) ke dalam metode bahasa pemrograman Java menggunakan anotasi khusus seperti \texttt{@Given}, \texttt{@When}, \texttt{@Then}, dan \texttt{@And}.
    \item \textbf{cucumber-junit-platform-engine}: Mesin penghubung yang menerjemahkan pengujian Cucumber agar dapat dibaca dan dieksekusi secara langsung oleh \textit{JUnit Platform Engine} (JUnit 5).
    \item \textbf{junit-platform-suite}: Menyediakan kapabilitas untuk mengelompokkan beberapa berkas pengujian ke dalam satu kesatuan uji (\textit{test suite}). Fitur ini sangat krusial dalam menyaring jalannya eksekusi berdasarkan kriteria tertentu seperti tag Gherkin.
\end{enumerate}

\section{Tugas 1: Implementasi Gherkin Syntax dan Cucumber}
Pada tugas pertama ini, diimplementasikan skenario-skenario pengujian login SauceDemo dengan memisahkan berkas skenario berbasis perilaku (Gherkin) dengan kode logika pengujian (Java) menggunakan pola desain \textit{Page Object Model} (POM).

\subsection{1. Skenario Gherkin (login.feature)}
Berkas \texttt{login.feature} mendeskripsikan perilaku login pengguna pada platform SauceDemo dengan berbagai variasi masukan data dan peran akun.

\begin{lstlisting}[caption=Berkas login.feature, language={}]
@Login
Feature: SauceDemo Login Scenarios

  Background:
    Given User is on the login page

  @StandardUser @Positive
  Scenario: Successful login with standard user
    When User enter username "standard_user" and password "secret_sauce"
    And User click the login button
    Then User should be redirected to the inventory page

  @LockedOut @Negative
  Scenario: Login failure with locked out user
    When User enter username "locked_out_user" and password "secret_sauce"
    And User click the login button
    Then User should see error message "Sorry, this user has been locked out."

  @ProblemUser @Positive
  Scenario: Login with problem user and verify broken image
    When User enter username "problem_user" and password "secret_sauce"
    And User click the login button
    Then User should be redirected to the inventory page
    And First inventory item image source should contain "sl-404"

  @PerformanceGlitch @Positive
  Scenario: Login with performance glitch user
    When User enter username "performance_glitch_user" and password "secret_sauce"
    And User click the login button
    Then User should be redirected to the inventory page with delay

  @ErrorUser @Positive
  Scenario: Login with error user
    When User enter username "error_user" and password "secret_sauce"
    And User click the login button
    Then User should be redirected to the inventory page

  @VisualUser @Positive
  Scenario: Login with visual user and verify inventory list is displayed
    When User enter username "visual_user" and password "secret_sauce"
    And User click the login button
    Then User should be redirected to the inventory page
    And Inventory list should be displayed
\end{lstlisting}

\subsection{Analisis Skenario Gherkin}
Komponen-komponen penyusun berkas fitur di atas dianalisis sebagai berikut:
\begin{itemize}
    \item \textbf{Tag Global (\texttt{@Login})}: Diletakkan pada baris pertama di atas definisi \textit{Feature}. Tag ini menandai bahwa seluruh skenario di dalam berkas ini merupakan bagian dari pengujian proses masuk (login).
    \item \textbf{Latar Belakang (\texttt{Background})}: Mendefinisikan langkah prasyarat awal (\texttt{Given User is on the login page}) yang akan dijalankan secara otomatis sebelum memulai setiap skenario pengujian individu. Hal ini mengeliminasi duplikasi penulisan langkah navigasi URL awal.
    \item \textbf{Skenario Uji (\texttt{Scenario})}: Membagi pengujian ke dalam alur tertentu berdasarkan kondisi akun pengguna.
    \item \textbf{Tag Khusus (\texttt{@StandardUser}, \texttt{@LockedOut}, \texttt{@Positive}, \texttt{@Negative}, dll.)}: Digunakan untuk melabeli karakteristik skenario. Skenario positif diberi tag \texttt{@Positive}, sedangkan skenario negatif (misal kegagalan login akibat akun terkunci) diberi tag \texttt{@Negative}. Tag ini juga dapat digunakan sebagai filter eksekusi.
    \item \textbf{Kata Kunci Gherkin}:
    \begin{itemize}
        \item \texttt{Given}: Menunjukkan kondisi awal atau prasyarat (dalam hal ini peramban mengarah ke halaman login).
        \item \texttt{When}: Menggambarkan aksi yang dilakukan oleh pengguna (mengisi data kredensial).
        \item \texttt{And}: Berfungsi sebagai langkah tambahan yang terhubung dengan langkah sebelumnya (menekan tombol masuk atau memverifikasi detail visual).
        \item \texttt{Then}: Menggambarkan hasil yang diharapkan setelah aksi dilakukan (berhasil masuk ke halaman produk atau menampilkan pesan kesalahan).
    \end{itemize}
\end{itemize}

\subsection{2. Page Object Model (POM) — basePage.java}
Kelas \texttt{basePage} bertindak sebagai kelas induk (\textit{super class}) yang menampung konfigurasi instansi WebDriver dan metode-metode pembantu umum untuk mengoperasikan elemen pada web.

\begin{lstlisting}[language=Java, caption=Kelas basePage.java]
package pages;

import org.openqa.selenium.By;
import org.openqa.selenium.WebDriver;
import org.openqa.selenium.WebElement;
import org.openqa.selenium.support.ui.ExpectedConditions;
import org.openqa.selenium.support.ui.WebDriverWait;

import java.time.Duration;

public class basePage {
    WebDriver driver;
    WebDriverWait wait;

    public basePage(WebDriver driver) {
        this.driver = driver;
        this.wait = new WebDriverWait(driver, Duration.ofSeconds(5));
    }

    public WebElement waitUntil(By by)
    {
        wait.until(ExpectedConditions.visibilityOfElementLocated(by));
        return driver.findElement(by);
    }

    public void click(By by)
    {
        waitUntil(by).click();
    }
    public void sendKeys(By by,String query)
    {
        waitUntil(by).sendKeys(query);
    }
}
\end{lstlisting}

\subsection{Analisis basePage}
Konstruktor kelas \texttt{basePage} menginisialisasi instansi \texttt{WebDriver} dan \texttt{WebDriverWait} dengan durasi batas tunggu eksplisit (\textit{explicit wait}) selama 5 detik. 
Metode \texttt{waitUntil()} memastikan bahwa program akan menunggu hingga elemen target yang didefinisikan oleh \texttt{By locator} terlihat secara visual di halaman web sebelum melakukan interaksi. Hal ini sangat berguna untuk meminimalisasi kegagalan pengujian akibat halaman yang lambat dimuat (\textit{flaky tests}). 
Metode \texttt{click()} dan \texttt{sendKeys()} memanfaatkan \texttt{waitUntil()} untuk melakukan klik dan pengiriman karakter (teks) ke elemen input secara aman.

\subsection{3. Page Object Model (POM) — loginPage.java}
Kelas \texttt{loginPage} mewarisi \texttt{basePage} dan berisi definisi lokasi elemen (locators) serta interaksi khusus yang terjadi di halaman login SauceDemo.

\begin{lstlisting}[language=Java, caption=Kelas loginPage.java]
package pages;

import org.openqa.selenium.By;
import org.openqa.selenium.WebDriver;

public class loginPage extends basePage {
    private WebDriver driver;

    public loginPage(WebDriver driver) {
        super(driver);
        this.driver = driver;
    }

    // Locators
    private final By usernameField = By.id("user-name");
    private final By passwordField = By.id("password");
    private final By loginButton = By.id("login-button");
    private final By errorMessage = By.cssSelector("[data-test='error']");

    // Actions
    public void enterUsername(String username) {
        sendKeys(usernameField, username);
    }

    public void enterPassword(String password) {
        sendKeys(passwordField, password);
    }

    public void clickLogin() {
        click(loginButton);
    }

    public String getErrorMessageText() {
        return waitUntil(errorMessage).getText();
    }
}
\end{lstlisting}

\subsection{Analisis loginPage}
Kelas \texttt{loginPage} mendefinisikan lokasi elemen input nama pengguna (\texttt{usernameField}), kata sandi (\texttt{passwordField}), tombol kirim (\texttt{loginButton}), dan elemen penampung teks galat (\texttt{errorMessage}) menggunakan pencarian berbasis ID dan CSS Selector. 
Metode-metode aksi seperti \texttt{enterUsername()}, \texttt{enterPassword()}, dan \texttt{clickLogin()} memanggil fungsi dasar \texttt{sendKeys()} dan \texttt{click()} yang diwarisi dari \texttt{basePage}, menjaga agar kelas ini tidak berurusan secara langsung dengan pemanggilan fungsi peramban yang rumit. Metode \texttt{getErrorMessageText()} mengembalikan pesan kesalahan ketika skenario masuk gagal dijalankan.

\subsection{4. Page Object Model (POM) — inventoryPage.java}
Kelas \texttt{inventoryPage} merepresentasikan halaman katalog produk setelah pengguna berhasil masuk.

\begin{lstlisting}[language=Java, caption=Kelas inventoryPage.java]
package pages;

import org.openqa.selenium.By;
import org.openqa.selenium.WebDriver;
import org.openqa.selenium.WebElement;

public class inventoryPage extends basePage {
    private WebDriver driver;

    public inventoryPage(WebDriver driver) {
        super(driver);
        this.driver = driver;
    }

    // Locators
    private final By firstItemImage = By.cssSelector(".inventory_item_img img");
    private final By inventoryList = By.className("inventory_list");

    // Actions & Validations
    public String getCurrentPageUrl() {
        return driver.getCurrentUrl();
    }

    public String getFirstItemImageSrc() {
        WebElement image = waitUntil(firstItemImage);
        return image.getAttribute("src");
    }

    public boolean isInventoryListDisplayed() {
        return waitUntil(inventoryList).isDisplayed();
    }
}
\end{lstlisting}

\subsection{Analisis inventoryPage}
Kelas \texttt{inventoryPage} digunakan untuk memvalidasi status halaman pasca-autentikasi. Elemen utama yang dipetakan adalah tautan gambar produk pertama (\texttt{firstItemImage}) dan container daftar produk (\texttt{inventoryList}). 
Metode \texttt{getCurrentPageUrl()} mengekstrak URL halaman aktif untuk memastikan pengalihan rute halaman berhasil ke arah \texttt{/inventory.html}. Metode \texttt{getFirstItemImageSrc()} digunakan khusus untuk mengidentifikasi apakah gambar produk rusak pada kasus pengguna bermasalah (\texttt{problem\_user}), sedangkan \texttt{isInventoryListDisplayed()} memastikan elemen daftar produk tampil di layar.

\subsection{5. Step Definitions dan Hooks — loginStepDef.java}
Kelas \texttt{loginStepDef} bertindak sebagai jembatan yang menghubungkan langkah-langkah deklaratif berkas Gherkin dengan logika eksekusi kode Java menggunakan objek-objek halaman (POM).

\begin{lstlisting}[language=Java, caption=Kelas loginStepDef.java]
package stepDef;

import io.cucumber.java.After;
import io.cucumber.java.Before;
import io.cucumber.java.en.And;
import io.cucumber.java.en.Given;
import io.cucumber.java.en.Then;
import io.cucumber.java.en.When;
import org.openqa.selenium.WebDriver;
import org.openqa.selenium.chrome.ChromeDriver;
import org.openqa.selenium.support.ui.ExpectedConditions;
import org.openqa.selenium.support.ui.WebDriverWait;
import pages.loginPage;
import pages.inventoryPage;

import java.time.Duration;

import static org.junit.jupiter.api.Assertions.assertEquals;
import static org.junit.jupiter.api.Assertions.assertTrue;

public class loginStepDef {
    private WebDriver driver;
    private loginPage login;
    private inventoryPage inventory;

    // Cucumber Hooks: @Before runs before each scenario
    @Before
    public void setUp() {
        driver = new ChromeDriver();
        driver.manage().window().maximize();
        login = new loginPage(driver);
        inventory = new inventoryPage(driver);
    }

    // Cucumber Hooks: @After runs after each scenario
    @After
    public void tearDown() {
        if (driver != null) {
            driver.quit();
        }
    }

    @Given("User is on the login page")
    public void userIsOnLoginPage() {
        driver.get("https://www.saucedemo.com/");
    }

    @When("User enter username {string} and password {string}")
    public void userEnterUsernameAndPassword(String username, String password) {
        login.enterUsername(username);
        login.enterPassword(password);
    }

    @And("User click the login button")
    public void userClickLoginButton() {
        login.clickLogin();
    }

    @Then("User should be redirected to the inventory page")
    public void userShouldBeRedirectedToInventoryPage() {
        String expectedUrl = "https://www.saucedemo.com/inventory.html";
        assertEquals(expectedUrl, inventory.getCurrentPageUrl());
    }

    @Then("User should see error message {string}")
    public void userShouldSeeErrorMessage(String expectedError) {
        String actualError = login.getErrorMessageText();
        assertTrue(actualError.contains(expectedError),
                "Error message should contain: " + expectedError + " but got: " + actualError);
    }

    @And("First inventory item image source should contain {string}")
    public void firstInventoryItemImageSourceShouldContain(String keyword) {
        String imageSrc = inventory.getFirstItemImageSrc();
        assertTrue(imageSrc != null && imageSrc.contains(keyword),
                "Image src should contain: " + keyword + " but got: " + imageSrc);
    }

    @Then("User should be redirected to the inventory page with delay")
    public void userShouldBeRedirectedToInventoryPageWithDelay() {
        WebDriverWait wait = new WebDriverWait(driver, Duration.ofSeconds(10));
        wait.until(ExpectedConditions.urlToBe("https://www.saucedemo.com/inventory.html"));
        assertEquals("https://www.saucedemo.com/inventory.html", inventory.getCurrentPageUrl());
    }

    @And("Inventory list should be displayed")
    public void inventoryListShouldBeDisplayed() {
        assertTrue(inventory.isInventoryListDisplayed(), "Inventory list should be displayed");
    }
}
\end{lstlisting}

\subsection{Analisis loginStepDef}
Metode-metode dalam kelas ini dianotasikan sesuai dengan langkah-langkah di berkas fitur:
\begin{itemize}
    \item \texttt{userEnterUsernameAndPassword(String username, String password)}: Menerima argumen berupa teks dari skenario Gherkin (menggunakan pola penangkap parameter \texttt{\{string\}}) lalu meneruskannya ke objek \texttt{loginPage}.
    \item \texttt{userShouldBeRedirectedToInventoryPage()}: Melakukan validasi pengalihan halaman dengan membandingkan URL saat ini menggunakan \texttt{assertEquals()}.
    \item \texttt{firstInventoryItemImageSourceShouldContain(String keyword)}: Memvalidasi apakah string URL gambar produk pertama mengandung kata kunci tertentu (seperti \texttt{"sl-404"}).
    \item \texttt{userShouldBeRedirectedToInventoryPageWithDelay()}: Menggunakan \texttt{WebDriverWait} lokal untuk memberikan toleransi waktu pemuatan halaman yang lebih panjang sebelum melakukan verifikasi pengalihan URL.
\end{itemize}

\subsection{6. Test Runner — JUnit 5 Suite (LoginSuiteRunner.java)}
Eksekusi berkas fitur Cucumber dilakukan melalui sebuah kelas perantara bernama \texttt{LoginSuiteRunner}.

\begin{lstlisting}[language=Java, caption=Kelas LoginSuiteRunner.java]
package runner;

import io.cucumber.junit.platform.engine.Constants;
import org.junit.platform.suite.api.ConfigurationParameter;
import org.junit.platform.suite.api.IncludeEngines;
import org.junit.platform.suite.api.SelectClasspathResource;
import org.junit.platform.suite.api.Suite;

@Suite
@IncludeEngines("cucumber")
@SelectClasspathResource("features/login.feature")
@ConfigurationParameter(key = Constants.GLUE_PROPERTY_NAME, value = "stepDef")
@ConfigurationParameter(key = Constants.FILTER_TAGS_PROPERTY_NAME, value = "@Login")
public class LoginSuiteRunner {
}
\end{lstlisting}

\subsection{Analisis LoginSuiteRunner}
Anotasi-anotasi penting pada kelas runner ini didefinisikan sebagai berikut:
\begin{itemize}
    \item \texttt{@Suite}: Mengonfigurasi kelas agar dapat dijalankan sebagai kumpulan pengujian terpadu JUnit Platform.
    \item \texttt{@IncludeEngines("cucumber")}: Menginstruksikan platform pengujian JUnit 5 untuk menggunakan mesin eksekusi Cucumber.
    \item \texttt{@SelectClasspathResource("features/login.feature")}: Menunjuk lokasi fisik dari berkas fitur Gherkin yang akan dieksekusi.
    \item \texttt{@ConfigurationParameter}: Menyediakan parameter konfigurasi bagi Cucumber.
    \begin{itemize}
        \item \texttt{Constants.GLUE\_PROPERTY\_NAME}: Menentukan nama paket Java yang menyimpan berkas pendefinisi langkah (\texttt{"stepDef"}).
        \item \texttt{Constants.FILTER\_TAGS\_PROPERTY\_NAME}: Menentukan ekspresi tag penyaring untuk memilih skenario tertentu. Pada konfigurasi awal Tugas 1, nilainya diset ke \texttt{"@Login"} untuk mengeksekusi semua skenario masuk di dalam berkas fitur.
    \end{itemize}
\end{itemize}

\subsection{7. Hasil Eksekusi Tugas 1}
Ketika kelas \texttt{LoginSuiteRunner} dijalankan dengan menyaring tag \texttt{"@Login"}, JUnit 5 Suite akan menjalankan seluruh 6 skenario login SauceDemo yang terdaftar pada berkas fitur. Peramban web akan terbuka secara bergantian sebanyak 6 kali untuk melakukan pengujian masing-masing peran akun.

% HINT GAMBAR: Masukkan screenshot hasil eksekusi 6 skenario login SauceDemo di IntelliJ IDEA
% \begin{center}
%     \includegraphics[width=15cm]{gambar/eksekusi_tugas1.png}
%     \captionof{figure}{Hasil Eksekusi 6 Skenario Login SauceDemo (@Login)}
% \end{center}

Berdasarkan hasil eksekusi tersebut, seluruh skenario uji berhasil mendapatkan status sukses (\textit{passed}) yang diindikasikan dengan simbol centang hijau di konsol IntelliJ IDEA. Hal ini membuktikan bahwa alur login untuk standard user, locked out user, problem user, performance glitch user, error user, dan visual user berhasil diproses sesuai kriteria masing-masing.

\section{Tugas 2: Eksplorasi Tag Gherkin dan Filter Tag pada Test Runner}
Pada tugas kedua, dilakukan analisis kapabilitas penyaringan pengujian menggunakan parameter penyaring pada Test Runner dengan target tag \texttt{@Negative}.

\subsection{Konfigurasi Filter Tag}
Konfigurasi penyaringan dilakukan dengan mengubah nilai properti filter tag pada kelas \texttt{LoginSuiteRunner} menjadi \texttt{"@Negative"}:

\begin{lstlisting}[language=Java, caption=Potongan Konfigurasi Filter Tag @Negative]
@ConfigurationParameter(
    key = Constants.FILTER_TAGS_PROPERTY_NAME, 
    value = "@Negative"
)
\end{lstlisting}

\subsection{Hasil Analisis \& Eksekusi}
Dengan mengganti nilai penyaring menjadi \texttt{"@Negative"}, JUnit Suite Runner akan memindai berkas skenario dan hanya meluncurkan skenario pengujian yang diberi label tag tersebut. 

Skenario yang berhasil dijalankan adalah skenario login gagal dengan akun terkunci (\texttt{locked\_out\_user}) karena merupakan satu-satunya skenario yang dilabeli dengan tag \texttt{@Negative}. Sementara itu, 5 skenario lainnya yang dilabeli dengan tag \texttt{@Positive} akan dilewati (\textit{skipped}).

% HINT GAMBAR: Masukkan screenshot hasil eksekusi filter tag @Negative
% \begin{center}
%     \includegraphics[width=15cm]{gambar/eksekusi_tugas2_passed.png}
%     \captionof{figure}{Hasil Eksekusi Filter Tag @Negative}
% \end{center}

% HINT GAMBAR: Masukkan screenshot konsol peringatan/keterangan filter mismatch di IntelliJ
% \begin{center}
%     \includegraphics[width=15cm]{gambar/eksekusi_tugas2_mismatch.png}
%     \captionof{figure}{Pesan Filter Mismatch pada Konsol Eksekusi}
% \end{center}

Konsol eksekusi JUnit 5 akan menampilkan informasi bahwa hanya satu tes yang berjalan sukses, sedangkan skenario lainnya tidak dieksekusi dikarenakan ketidakcocokan tag (\textit{filter mismatch}). Hal ini menunjukkan efektivitas pemisahan uji dalam mempercepat verifikasi skenario kegagalan sistem.

\section{Tugas 3: Penerapan Saringan Test Spesifik via Test Runner}
Tugas ketiga menguji penerapan penyaringan yang lebih spesifik untuk menjalankan skenario login khusus pengguna standar (\texttt{@StandardUser}).

\subsection{Konfigurasi Filter}
Berkas \texttt{LoginSuiteRunner.java} disesuaikan kembali dengan mengubah nilai parameter \texttt{FILTER\_TAGS\_PROPERTY\_NAME} menjadi \texttt{"@StandardUser"}:

\begin{lstlisting}[language=Java, caption=Potongan Konfigurasi Filter Tag @StandardUser]
@ConfigurationParameter(
    key = Constants.FILTER_TAGS_PROPERTY_NAME, 
    value = "@StandardUser"
)
\end{lstlisting}

\subsection{Hasil Analisis \& Eksekusi}
Saat Test Suite dipicu, sistem hanya mendeteksi satu kecocokan skenario pengujian, yaitu masuk sukses dengan akun standar (\texttt{standard\_user}). Sistem secara otomatis mengabaikan skenario login negatif maupun variasi pengguna positif lainnya.

% HINT GAMBAR: Masukkan screenshot hasil eksekusi filter tag @StandardUser
% \begin{center}
%     \includegraphics[width=15cm]{gambar/eksekusi_tugas3_passed.png}
%     \captionof{figure}{Hasil Eksekusi Filter Tag @StandardUser}
% \end{center}

% HINT GAMBAR: Masukkan screenshot konsol filter mismatch non-StandardUser
% \begin{center}
%     \includegraphics[width=15cm]{gambar/eksekusi_tugas3_mismatch.png}
%     \captionof{figure}{Pesan Filter Mismatch Skenario Non-StandardUser}
% \end{center}

Hasil konsol menunjukkan satu skenario sukses dijalankan dan lima skenario lainnya ditandai sebagai diabaikan. Fitur filter tag spesifik ini mempermudah proses debugging ketika terjadi kesalahan logika hanya pada bagian alur tertentu tanpa perlu membuang waktu menjalankan seluruh suite pengujian.

\section{Poin Tambahan (Bonus): Penerapan Hooks}
Dalam implementasi BDD Cucumber, pengelolaan siklus hidup pengujian dikelola secara dinamis menggunakan anotasi Hooks (\texttt{@Before} dan \texttt{@After}).

\subsection{Implementasi Hooks}
Implementasi kode program Hooks diletakkan pada kelas pendefinisi langkah (\texttt{loginStepDef.java}):

\begin{lstlisting}[language=Java, caption=Implementasi Hooks pada loginStepDef.java]
@Before
public void setUp() {
    driver = new ChromeDriver();
    driver.manage().window().maximize();
    login = new loginPage(driver);
    inventory = new inventoryPage(driver);
}

@After
public void tearDown() {
    if (driver != null) {
        driver.quit();
    }
}
\end{lstlisting}

\subsection{Analisis Peran dan Keuntungan Hooks}
Penerapan Cucumber Hooks memberikan beberapa keuntungan struktural bagi kode otomatisasi pengujian:
\begin{enumerate}
    \item \textbf{Manajemen Siklus Peramban (Browser Lifecycle Management)}: Dengan anotasi \texttt{@Before}, instansi Chrome WebDriver yang baru dan segar akan disiapkan sebelum memulai langkah pertama dari setiap skenario pengujian. Hal ini memastikan tidak ada sisa-sisa \textit{cookies} atau \textit{session} dari pengujian sebelumnya yang dapat mengganggu validitas hasil uji.
    \item \textbf{Pembersihan Otomatis Sumber Daya (Automatic Resource Cleanup)}: Anotasi \texttt{@After} menjamin peramban web akan ditutup menggunakan metode \texttt{driver.quit()} secara otomatis setelah pengujian selesai, baik pengujian tersebut berakhir sukses maupun gagal. Hal ini mencegah terjadinya kebocoran memori (\textit{memory leak}) pada sistem operasi akibat proses peramban yang menggantung.
    \item \textbf{Reusabilitas Kode (Code Reusability)}: Kode pembuka dan penutup peramban dikumpulkan dalam satu tempat sentral sehingga tidak perlu didefinisikan secara repetitif di setiap langkah pengujian individu. Hal ini sejalan dengan prinsip desain kode bersih (\textit{Clean Code}).
\end{enumerate}

\section{Kode Sumber Lengkap di GitHub}
Seluruh kode sumber program pengujian otomatis berbasis BDD Gherkin dan Cucumber ini telah diunggah ke dalam repositori GitHub publik melalui tautan berikut:
\begin{center}
    \url{https://github.com/akmalmanggala/BDD-gherkin.git}
\end{center}


% BAB 3: KESIMPULAN
\chapter{Kesimpulan}

\section{Kesimpulan}
Berdasarkan hasil praktikum yang telah diselesaikan mengenai otomatisasi pengujian perangkat lunak menggunakan metode \textit{Behavior Driven Development} (BDD) dengan Cucumber, diperoleh kesimpulan sebagai berikut:

\begin{enumerate}
    \item Metodologi BDD melalui penulisan sintaks Gherkin terbukti mempermudah komunikasi dan penyelarasan pemahaman alur bisnis antara tim non-teknis dengan tim pengembang karena deskripsi pengujian ditulis dalam bahasa natural terstruktur.
    \item Pola desain \textit{Page Object Model} (POM) dengan kelas induk \texttt{basePage} dapat berkolaborasi dengan baik bersama Cucumber, menghasilkan pemisahan yang bersih antara deklarasi alur perilaku pengguna (di berkas fitur), penafsiran langkah aksi (di kelas step definitions), dan definisi interaksi elemen halaman web (di kelas page objects).
    \item Penggunaan Cucumber Hooks (\texttt{@Before} dan \texttt{@After}) memberikan jaminan pemeliharaan siklus hidup peramban secara otomatis dan konsisten, meminimalisasi overhead pemakaian memori, serta meningkatkan reusabilitas kode program.
    \item Fitur filter tag di JUnit 5 Suite Runner terbukti sangat efektif untuk mengelompokkan dan memilih skenario pengujian spesifik yang ingin dijalankan. Hal ini mempercepat waktu umpan balik pengujian selama proses pengembangan perangkat lunak berlangsung.
\end{enumerate}

Secara keseluruhan, integrasi JUnit 5 Platform Suite, Cucumber, dan Selenium WebDriver menghasilkan solusi pengujian otomatis yang kokoh, modular, mudah dirawat, dan siap untuk diintegrasikan dalam siklus rilis berkelanjutan (\textit{Continuous Integration}).

\end{document}
